\cotefile{136}
\scanpage{111}
\begin{center}Foncteurs de De Rham généralisés\end{center}
On s'intéresse aux foncteurs
\[
 C^{**}:(\mathrm{Ss})^{o}\longrightarrow
 \left[\begin{aligned}
 &S_k\text{-algèbres bigraduées \inserted{alternées}}\\
 &\text{différentielles homogènes}\\
 &\text{\unsure{ties|tées} à p.d. (bidegrés}\\
 &\text{positifs, diff. de degré}\\
 &+1,\ \ldots\ldots\,)
 \end{aligned}\right]=\unsure{\mathscr S_k\mathscr D|S_k\mathscr D|\mathfrak D}
\]
\margin{$\unsure{o|\circ}$}
et on s'intéresse à ceux qui transforment $\underrightarrow{\lim}$ quelconques de $\mathrm{Ss}$ en $\underleftarrow{\lim}$, donc qui sont représentables.
Ils \deleted{\unsure{doivent|donnent}} \inserted{peuvent} donc s'interpréter \inserted{(\unsure{fonctoriellement|ponctuellement|fantasmes})} comme des $S_k$-algèbres bigraduées alternées \ldots\ $1/2$ simpliciales, i.e. des foncteurs
\[
 \mathbb{\Delta}^{o}\longrightarrow\mathscr D
\]
où on identifie dans les notations les deux sortes d'objets, écrivant \unsure{en|un}
\[
 C^{**}=C_*^{**}=(C_n^{**})_{n\geq 0}\quad\ldots
\]
et de $=$ pour les morphismes entre foncteurs, correspondant aux morphismes entre ensembles $1/2$ simpliciaux, respectant toutes les structures.

À $C_*^{**}$ est associé un foncteur ``cohomologie de De Rham associée à $C_{\unsure{*}}^{**}$''
\[
 \mathscr H_C^{**}(\mathscr X)=H_{\unsure{\mathscr S|\mathrm{gr}|\mathrm{G}|\varepsilon}}^{**}(C^{**}(\mathscr X))
\]
à valeurs dans les $S_k$-algèbres bigraduées alternées. On s'intéresse aux $C^{**}\deleted{(\mathscr X)}$ (dits admissibles) pour lesquels, pour tt $\mathscr X$,
% archive sheet number: 110
